\subsection{Komparatorschaltung zur Akku-Überwachung berechnen}
% Alt: Aufgabe 10
\Aufgabe{
    \label{Aufg10}
    Zur Überwachung des Ladezustandes eines Akkus soll ein Komparator eingesetzt werden. Der Komparator soll umschalten, wenn eine Batteriespannung
    $U_\mathrm{bat}$ von $6\,\mathrm{V}$ unterschritten wird. Der Widerstand $R_1$ beträgt $10\,\mathrm{k\Omega}$. Die Betriebsspannung des Komparators beträgt $\pm 15\,\mathrm{V}$.
    \begin{itemize}
        \item [a)] Wie groß muss der Widerstand $R_2$ gewählt werden?
        \item [b)] Zur Anzeige des Ladezustands soll eine LED verwendet werden. Erweitern Sie die Schaltung um die Ladungsanzeige per LED.
        \item [c)] Durch die LED mit der Flussspannung von 2,2 V darf maximal ein Strom von 20 mA fließen. Wie groß muss der Vorwiderstand gewählt werden?
        \item [d)] Die Spannung am Spannungsteiler ist abhängig von der Betriebsspannung. Welche anderen geeigneten Bauelemente existieren für die Erzeugung der Referenzspannung? Nennen und zeichnen Sie ein Beispiel.
    \end{itemize}
    \begin{figure}[H]
        \centering
        % Linkes Bild
        \begin{subfigure}[b]{0.45\textwidth}
            \centering
            \input{Tikz/src/FigAkkuUeberwachung.tex}\alttext{Das Bild zeigt zwei Teile: oben ein Schaltplan und unten eine zugehörige Grafik.
                Im Schaltplan ist eine Spannungsquelle mit +15 Volt angeschlossen an einen Spannungsteiler bestehend aus zwei Widerständen \(R_1\) und \(R_2\), die in Serie geschaltet sind. Der Verbindungspunkt zwischen \(R_1\) und \(R_2\) ist mit dem nicht-invertierenden Eingang (+) eines Operationsverstärkers verbunden. Der invertierende Eingang (-) des Operationsverstärkers ist mit einer Batterie \(U_{\mathrm{bat}}\) verbunden, deren negativer Pol auf Masse liegt. Der Ausgang des Operationsverstärkers liefert die Ausgangsspannung \(U_\mathrm{A}\).}
        \end{subfigure}

        \begin{subfigure}[b]{0.45\textwidth}
            \centering
            \includesvg[width=\textwidth]{Bilder/Aufgaben/FigAkkuUeberwachung}\alttext{ Die darunterliegende Grafik zeigt den Zusammenhang zwischen der Eingangsspannung \(U_\mathrm{E}\) (x-Achse) und der Ausgangsspannung \(U_\mathrm{A}\) (y-Achse). Die Kurve verläuft als Rechteckfunktion: Für positive Werte von \(U_\mathrm{E}\) größer als eine Referenzspannung \(U_{\mathrm{Ref}}\) ist \(U_\mathrm{A}\) konstant auf \(+U_\mathrm{B}\); für negative Werte von \(U_\mathrm{E}\) kleiner als \(-U_{\mathrm{Ref}}\) ist \(U_\mathrm{A}\) konstant auf \(-U_\mathrm{B}\). Zwischen diesen Bereichen liegt eine scharfe Umschaltung bei \(±U_{\mathrm{Ref}}\).}
        \end{subfigure}
        \label{fig:FigAkkuUeberwachung}
    \end{figure}
}

\Loesung{
    \begin{itemize}
        \item[a)] Berechnung des Widerstands $R_2$

              Der Komparator soll umschalten, wenn die Batteriespannung $U_{\text{bat}}$ unter $6\,\mathrm{V}$ fällt. Das Verhältnis der Widerstände $R_1$ und $R_2$ im Spannungsteiler muss so gewählt werden, dass am nicht-invertierenden Eingang die Referenzspannung $U_{\text{ref}} = 6\,\mathrm{V}$ anliegt:
              \[
                  U_{\text{ref}} = \frac{R_2}{R_1 + R_2} \cdot U_\mathrm{B}
              \]
              mit $U_\mathrm{B} = 15\,\mathrm{V}$ und $R_1 = 10\,\mathrm{k}\Omega$.
              Einsetzen:
              \[
                  6\,\mathrm{V} = \frac{R_2}{10\,\mathrm{k}\Omega + R_2} \cdot 15\,\mathrm{V}
              \]
              Umstellen nach $R_2$:
              \[
                  R_2 = \frac{6\,\mathrm{V} \cdot 10\,\mathrm{k}\Omega}{15\,\mathrm{V} - 6\,\mathrm{V}} = 6,66\,\mathrm{k}\Omega
              \]

        \item[b)] Erweiterung der Schaltung um eine LED-Anzeige:

              Die LED soll anzeigen, wenn die Batteriespannung unter $6\,\mathrm{V}$ fällt. Dazu wird die LED mit Vorwiderstand $R_V$ an den Ausgang des Komparators geschaltet.

        \item Schaltverhalten der Schaltung:

              \begin{itemize}
                  \item Leere Batterie ($U_{\text{bat}} = 5\,\mathrm{V}$):
                        \begin{align*}
                            U_+ & = 6\,\mathrm{V} - U_{\text{bat}}                \\
                            U_+ & = 6\,\mathrm{V} - 5\,\mathrm{V} = 1\,\mathrm{V}
                        \end{align*}
                        Da $U_+ < U_{\text{ref}}$, schaltet der Komparator um:
                        \[
                            U_\mathrm{A} = -15\,\mathrm{V}
                        \]
                        Die LED leuchtet.

                  \item Volle Batterie ($U_{\text{bat}} = 7\,\mathrm{V}$):
                        \begin{align*}
                            U_+ & = 6\,\mathrm{V} - U_{\text{bat}}                 \\
                            U_+ & = 6\,\mathrm{V} - 7\,\mathrm{V} = -1\,\mathrm{V}
                        \end{align*}
                        Da $U_+ > U_{\text{ref}}$, schaltet der Komparator:
                        \[
                            U_\mathrm{A} = 15\,\mathrm{V}
                        \]
                        Die LED bleibt aus.
              \end{itemize}

        \item[c)] Dimensionierung des Vorwiderstands $R_\mathrm{V}$

              Gegeben:
              \begin{itemize}
                  \item Flussspannung der LED: $U_\mathrm{F} = 2,2\,\mathrm{V}$
                  \item Strom durch die LED: $I_\mathrm{F} = 20\,\mathrm{mA}$
                  \item Ausgangsspannung des Komparators: $U_\mathrm{A} = -15\,\mathrm{V}$
              \end{itemize}

              \begin{figure}[H]
                  \centering
                  \begin{tikzpicture}
    \ctikzset{tripoles/en amp/input height=-0.45}
    \draw (0,0)node[en amp](E){};
    \draw (E.out) to[short, -o] (2,0)to[o-, R, l=$R_\mathrm{V}$, v<, name=RV] (6,0);
    \draw (6,-2)node[ground]{} to[led, i>, v, name=D] (6,0);
    \draw (E.+) to[short, -*] ( -4,0.5);

    \draw (E.-) -- (-2,-0.5) to[battery2, name=Ub, -*] (-2,-2)node[ground]{};
    \draw (2,-2) 
        to[short, o-] (-4,-2) 
        to[R,l=$R_\mathrm{2}$] (-4,0.5) 
        to[R, l=$R_\mathrm{1}$, -o] (-4,3) node[above]{$+15 \, \mathrm{V}$};
    
    \draw (2,0) to[open, v, name=ue] (2,-2);
    \draw (-3,0.5) to[open, v, name=ux] (-3,-2);
    \draw (-5,2) to[open, v, name=ua] (-5,-2);

    \varrmore{ue}{$U_\mathrm{A}$};
    \varrmore{Ub}{$U_\mathrm{bat}$};
    \iarrmore{D}{$I_\mathrm{F}$};
    \varrmore{D}{$U_\mathrm{F}$};
    \varrmore{RV}{$U_\mathrm{RV}$};
\end{tikzpicture}
                  \alttext{Die Abbildung zeigt eine Schaltung mit einem Operationsverstärker, der mit einer positiven Versorgungsspannung von \(+15 \ \mathrm{V}\) betrieben wird. Der Operationsverstärker ist als ideales Bauelement dargestellt und besitzt einen nichtinvertierenden Eingang (+), einen invertierenden Eingang (−) sowie einen Ausgang.
                      Der nichtinvertierende Eingang des Operationsverstärkers ist mit einem Spannungsteiler verbunden. Dieser besteht aus den Widerständen \(R_1\) und \(R_2\), die zwischen \(+15 \ \mathrm{V}\) und Masse geschaltet sind.
                      Der invertierende Eingang ist über eine Spannungsquelle \(U_{\mathrm{bat}}\) mit Masse verbunden.
                      Am Ausgang des Operationsverstärkers befindet sich ein Vorwiderstand \(R_\mathrm{V}\) Hinter diesem Widerstand ist eine LED angeschlossen, die gegen Masse geschaltet ist. Für die LED sind der Durchlassstrom \(I_\mathrm{F}\) sowie die Durchlassspannung \(U_\mathrm{F}\) eingezeichnet. Über dem Widerstand \(R_\mathrm{V}\) ist die Spannung \(U_{\mathrm{RV}}\) eingezeichnet.
                      Zusätzlich ist zwischen Ausgang und Masse die Ausgangsspannung \(U_\mathrm{A}\) eingezeichnet.}
                  \label{fig:LsgAkkuUeberwachung1}
              \end{figure}

              Der Vorwiderstand $R_\mathrm{V}$ berechnet sich aus:
              \[
                  R_\mathrm{V} = \frac{-U_\mathrm{A} - U_\mathrm{F}}{I_\mathrm{F}}
              \]
              Einsetzen der Werte:
              \[
                  R_\mathrm{V} = \frac{-(-15\,\mathrm{V}) - 2,2\,\mathrm{V}}{20\,\mathrm{mA}} = \frac{12,8\,\mathrm{V}}{0,02\,\mathrm{A}} = 640\,\Omega
              \]

              Für einen sicheren Betrieb der LED wird Widerstandswert von mindestens $R_\mathrm{V} = 640\,\Omega$ benötigt.
        \item[d)] Alternative Spannungsreferenz:

              Da die Spannung des Spannungsteilers abhängig von der Betriebsspannung ist, könnte eine Zener-Diode als Referenzelement verwendet werden. Eine Zener-Diode stellt eine konstante Referenzspannung unabhängig von Schwankungen der Versorgungsspannung bereit.

        \item Beispiel-Schaltung:
              \begin{itemize}
                  \item Die Zener-Diode wird mit einem Vorwiderstand an die Betriebsspannung $+15\,\mathrm{V}$ angeschlossen.
                  \item Die Referenzspannung wird direkt von der Kathode der Zener-Diode abgegriffen.
              \end{itemize}
    \end{itemize}
    \begin{figure}[H]
        \centering
        \input{Tikz/src/LsgAkkuUeberwachung4.tex}
        \alttext{\alttext{Die Abbildung zeigt eine Schaltung mit einem Operationsverstärker, der mit einer positiven Versorgungsspannung von \(+15 \ \mathrm{V}\) betrieben wird. Der Operationsverstärker ist als ideales Bauelement dargestellt und besitzt einen nichtinvertierenden Eingang (+), einen invertierenden Eingang (-) sowie einen Ausgang.
                Der nichtinvertierende Eingang ist mit einem Knoten verbunden, der über einen Widerstand von \(10 \ \mathrm{k \Ohm}\) an \(+15\,\mathrm{V}\) liegt. Zwischen diesem Knoten und Masse befindet sich eine Z-Diode mit einer angegebenen Zenerspannung von \(6 \ \mathrm{V}\) Dadurch entsteht am nichtinvertierenden Eingang eine stabilisierte Referenzspannung von etwa \(6 \ \mathrm{V}\). Der durch die Z-Diode fließende Strom ist mit \(I\) gekennzeichnet.
                Der invertierende Eingang ist über eine Spannungsquelle \(U_{\mathrm{bat}}\) mit Masse verbunden. Diese Spannungsquelle stellt die zu vergleichende Eingangsspannung dar.
                Am Ausgang des Operationsverstärkers ist ein Widerstand von \(640 \ \Omega\) angeschlossen. Zwischen diesem Widerstand und dem Bezugspunkt ist eine LED geschaltet. Die LED ist so eingebunden, dass sie abhängig vom Ausgangspotential des Operationsverstärkers Strom führt und damit den Schaltzustand optisch anzeigt.#
                Zusätzlich ist die Spannung \(U_\mathrm{a}\) zwischen dem Referenzknoten und Masse eingezeichnet.}}
        \label{fig:LsgAkkuUeberwachung4}
    \end{figure}
    Siehe: Abschnitt \ref{fig: Verschaltung Verstaerker} und Tabelle \ref{tab:Grundschaltungen}

}